\documentclass[10pt,a4paper]{amsart}
\usepackage{amsmath,amssymb,amsthm,graphicx}
% Retained minimal preamble for the Stacks Project chapter more-algebra.tex (More on Algebra).
% The chapter file pinned-original-source.tex carries only \input{preamble} on its first line, so no
% preamble text lives inside the chapter itself.  The Stacks Project uses one single shared preamble for
% every chapter; the snapshot stages it identically as sources/stacks-chow/preamble.tex and
% sources/stacks-etale-cohomology/preamble.tex.  The block below is copied mechanically, line for line,
% from sources/stacks-chow/preamble.tex lines 14-18 (the Xy-pic setup) and 154-217 (the native theorem
% declarations together with the mathematical macros).  Nothing was added, removed or redefined; the
% wrapper class and the amsmath/amssymb/amsthm/graphicx base packages are supplied by the assembler.
% Chapter: more-algebra.tex -- More on Algebra.
% For commutative diagrams we use Xy-pic
\usepackage[all]{xy}

% We use 2cell for 2-commutative diagrams.
\xyoption{2cell}

% Theorem environments.
%
\theoremstyle{plain}
\newtheorem{theorem}[subsection]{Theorem}
\newtheorem{proposition}[subsection]{Proposition}
\newtheorem{lemma}[subsection]{Lemma}

\theoremstyle{definition}
\newtheorem{definition}[subsection]{Definition}
\newtheorem{example}[subsection]{Example}
\newtheorem{exercise}[subsection]{Exercise}
\newtheorem{situation}[subsection]{Situation}

\theoremstyle{remark}
\newtheorem{remark}[subsection]{Remark}
\newtheorem{remarks}[subsection]{Remarks}

\numberwithin{equation}{subsection}

% Macros
%
\def\lim{\mathop{\mathrm{lim}}\nolimits}
\def\colim{\mathop{\mathrm{colim}}\nolimits}
\def\Spec{\mathop{\mathrm{Spec}}}
\def\Hom{\mathop{\mathrm{Hom}}\nolimits}
\def\Ext{\mathop{\mathrm{Ext}}\nolimits}
\def\SheafHom{\mathop{\mathcal{H}\!\mathit{om}}\nolimits}
\def\SheafExt{\mathop{\mathcal{E}\!\mathit{xt}}\nolimits}
\def\Sch{\mathit{Sch}}
\def\Mor{\mathop{\mathrm{Mor}}\nolimits}
\def\Ob{\mathop{\mathrm{Ob}}\nolimits}
\def\Sh{\mathop{\mathit{Sh}}\nolimits}
\def\NL{\mathop{N\!L}\nolimits}
\def\CH{\mathop{\mathrm{CH}}\nolimits}
\def\proetale{{pro\text{-}\acute{e}tale}}
\def\etale{{\acute{e}tale}}
\def\QCoh{\mathit{QCoh}}
\def\Ker{\mathop{\mathrm{Ker}}}
\def\Im{\mathop{\mathrm{Im}}}
\def\Coker{\mathop{\mathrm{Coker}}}
\def\Coim{\mathop{\mathrm{Coim}}}

% Boxtimes
%
\DeclareMathSymbol{\boxtimes}{\mathbin}{AMSa}{"02}

%
% Macros for moduli stacks/spaces
%
\def\QCohstack{\mathcal{QC}\!\mathit{oh}}
\def\Cohstack{\mathcal{C}\!\mathit{oh}}
\def\Spacesstack{\mathcal{S}\!\mathit{paces}}
\def\Quotfunctor{\mathrm{Quot}}
\def\Hilbfunctor{\mathrm{Hilb}}
\def\Curvesstack{\mathcal{C}\!\mathit{urves}}
\def\Polarizedstack{\mathcal{P}\!\mathit{olarized}}
\def\Complexesstack{\mathcal{C}\!\mathit{omplexes}}
% \Pic is the operator that assigns to X its picard group, usage \Pic(X)
% \Picardstack_{X/B} denotes the Picard stack of X over B
% \Picardfunctor_{X/B} denotes the Picard functor of X over B
\def\Pic{\mathop{\mathrm{Pic}}\nolimits}
\def\Picardstack{\mathcal{P}\!\mathit{ic}}
\def\Picardfunctor{\mathrm{Pic}}
\def\Deformationcategory{\mathcal{D}\!\mathit{ef}}

% --- appended 2026-10-12 (Chief): the project own definitions of the reference/slogan/history
% display environments, copied mechanically from the shared Stacks preamble lines 8-12
% (\usepackage{verbatim} plus the three \newenvironment lines, which the chapter preamble keeps
% verbatim from the project preamble). Nothing added, removed or redefined: these environments
% carry no mathematical content and the project itself typesets them as comments. Without these
% three lines a counted statement whose own text uses them fails with Environment slogan undefined.
\usepackage{verbatim}
\newenvironment{reference}{\comment}{\endcomment}
\newenvironment{slogan}{\comment}{\endcomment}
\newenvironment{history}{\comment}{\endcomment}
\usepackage{cleveref}
\crefname{originallemma}{original lemma}{original lemmas}
\Crefname{originallemma}{Original lemma}{Original lemmas}
\makeatletter
\newlabel{more-algebra-lemma-make-nonzero-divisor}{{068P}{source}{}{sourcepointer-1}{}}
\newlabel{more-algebra-lemma-make-nonzero-divisor@cref}{{[originallemma][0][]068P}{[1][0][]source}{}{}{}}
\makeatother
\crefname{originallemma}{original lemma}{original lemmas}
\Crefname{originallemma}{Original lemma}{Original lemmas}
\makeatletter
\newlabel{algebra-lemma-zero-at-ass-zero}{{0311}{source}{}{sourcepointer-2}{}}
\newlabel{algebra-lemma-zero-at-ass-zero@cref}{{[originallemma][0][]0311}{[1][0][]source}{}{}{}}
\makeatother
\crefname{originalreference}{original reference}{original references}
\Crefname{originalreference}{Original reference}{Original references}
\makeatletter
\newlabel{McCoy}{{cite-McCoy}{source}{}{sourcepointer-3}{}}
\newlabel{McCoy@cref}{{[originalreference][0][]cite-McCoy}{[1][0][]source}{}{}{}}
\makeatother
\title{A linear combination of elements defining an injective map is a nonzerodivisor in the polynomial ring}
\author{The Stacks Project Authors}
\date{stacks-project:more-algebra.tex:local-snapshot-20261012; GNU FDL 1.2 or later}
\begin{document}
\maketitle
\setcounter{subsection}{1}
\noindent\textit{Editorial note.} All delivered mathematical text is the author's own, copied from the pinned Stacks Project chapter file more-algebra.tex; no statement or proof step was rewritten, reordered or removed.\par\medskip
\noindent\textit{Editorial note.} The chapter file pinned-original-source.tex carries only \textbackslash{}input\{preamble\} on its first line, so no preamble text lives inside the chapter. The retained file sources/stacks-more-algebra/retained-preamble.tex is a mechanical, line-for-line copy of the project's single shared preamble as staged in the snapshot (Xy-pic setup plus theorem declarations and mathematical macros). Nothing was added, removed or redefined, and the excerpt serves only as preamble, never as a span coordinate.\par\medskip
\noindent\textit{Editorial note.} Both counted spans start at the \textbackslash{}begin line of their environment, so the statement and the author's proof are delivered whole; surrounding prose of the section is deliberately outside the counted spans and is carried as standard prerequisites instead.\par\medskip
\noindent\textit{Editorial note.} The citation fields (source\_id, source\_version, authors, source\_url, license.name, source\_title) are identical to those of every other card this lane has produced from this chapter, so the batch selector treats them as one source.\par\medskip
\noindent\textit{Editorial note.} The statement contains a comment-like 'reference' environment whose content cites McCoy. In this chapter those environments are defined with comment/endcomment, so the block is not typeset, the citation is never processed and no bibliography entry is required; the retained preamble is taken in full through its current last line so that the environment definition is present.\par\medskip
\noindent\textit{Allowed standard prerequisites (editorial).}\begin{itemize}
\item Associated primes of a module over a Noetherian ring and the criterion that a map of finite modules is injective if it is injective after localising at every associated prime. the author cites Algebra, Lemma algebra-lemma-zero-at-ass-zero at proof lines 7787-7789. Supplies the localisation criterion that turns the verification at associated primes into injectivity.
\item Symmetric powers of a free module are free, and the multiplication by a linear form on them. used at proof lines 7772-7786. Provides the maps whose injectivity is checked after localisation.
\item Every ring is a filtered union of Noetherian subrings, and injectivity is preserved by such unions. used at proof lines 7791-7795. Reduces the general case to the Noetherian one.
\end{itemize}\medskip
\ifdefined\hypertarget\hypertarget{sourcepointer-1}{}\fi
\noindent\textit{Source reference.} Original lemma 068P: focused-seeds/AGP/sources/stacks-snapshot-20261011/tags.txt, source lines 8013--8013. The lemma's own canonical tag, quoted so that readers can resolve the label inside the delivered text; it points at the counted statement itself and carries no proof content.\par\medskip
\ifdefined\hypertarget\hypertarget{sourcepointer-2}{}\fi
\noindent\textit{Source reference.} Original lemma 0311: focused-seeds/AGP/sources/stacks-snapshot-20261011/tags.txt, source lines 3968--3968. Navigation pointer to a statement the counted proof refers to but does not contain; nothing load-bearing is delegated to it.\par\medskip
\ifdefined\hypertarget\hypertarget{sourcepointer-3}{}\fi
\noindent\textit{Source reference.} Original reference cite-McCoy: sources/stacks-more-algebra/pinned-original-source.tex, source lines 7757--7757. Bibliographic pointer appearing in the counted text; it carries no mathematical argument.\par\medskip
\begin{lemma}
\label{lemma-make-nonzero-divisor}
\begin{reference}
This is a particular case of \cite[Corollary]{McCoy}
\end{reference}
Let $R$ be a ring. Let $a_1, \ldots, a_n \in R$ be elements such
that $R \to R^{\oplus n}$, $x \mapsto (xa_1, \ldots, xa_n)$ is injective.
Then the element $\sum a_i t_i$ of the polynomial ring $R[t_1, \ldots, t_n]$
is a nonzerodivisor.
\end{lemma}

\begin{proof}
If one of the $a_i$ is a unit this is just the statement that any
element of the form $t_1 + a_2 t_2 + \ldots + a_n t_n$ is a nonzerodivisor
in the polynomial ring over $R$.

\medskip\noindent
Case I: $R$ is Noetherian. Let $\mathfrak q_j$, $j = 1, \ldots, m$
be the associated primes of $R$. We have to show that
each of the maps
$$
\sum a_i t_i :
\text{Sym}^d(R^{\oplus n})
\longrightarrow
\text{Sym}^{d + 1}(R^{\oplus n})
$$
is injective. As $\text{Sym}^d(R^{\oplus n})$ is a free $R$-module its
associated primes are $\mathfrak q_j$, $j = 1, \ldots, m$. For each $j$
there exists an $i = i(j)$ such that $a_i \not \in \mathfrak q_j$ because
there exists an $x \in R$ with $\mathfrak q_jx = 0$ but $a_i x \not = 0$
for some $i$ by assumption. Hence $a_i$ is a unit in $R_{\mathfrak q_j}$
and the map is injective after localizing at $\mathfrak q_j$. Thus the map
is injective, see
Algebra, Lemma \ref{algebra-lemma-zero-at-ass-zero}.

\medskip\noindent
Case II: $R$ general. We can write $R$ as the union of Noetherian
rings $R_\lambda$ with $a_1, \ldots, a_n \in R_\lambda$. For each $R_\lambda$
the result holds, hence the result holds for $R$.
\end{proof}

\end{document}
